\scanpage{88}
% archive foliation: 87
\begin{center}- 139 -\end{center}
\hfill n\textsuperscript{o} 316

a) Soit $M'_s$ le sous-espace de $M_s$ engendré par les éléments de la forme $\beta_1\otimes\beta_2\otimes\cdots\otimes\beta_s$, où $\beta_i\in\mathcal L(V)$ et où $\beta_i=1$ pour au moins un $i$. Soient $z$ un élément de $U_s$, $z_0=\sum_{1\leq i_1,\ldots,i_s\leq m}a_{i_1\ldots i_s}\alpha_{i_1}\cdots\alpha_{i_s}$ sa composante dans $U^s$, où les $a_{i_1\ldots i_s}$ sont symétriques par rapport aux permutations des indices. Montrer que
\[
\pi_s(z)\equiv s!\sum_{1\leq i_1,\ldots,i_s\leq m}a_{i_1\ldots i_s}\alpha_{i_1}\otimes\cdots\otimes\alpha_{i_s}\quad(\mathrm{mod.}\ M'_s).
\]

b) Déduire de a), que, pour tout $z\in U$, il existe une représentation de dimension finie $\pi$ de $U$ telle que $\pi(z)\neq0$. (Pour une suite de cet exerc., cf.\ §7, exerc.\ 3).

6) - On suppose que $K$ est un corps. Soient $g$ une $K$-algèbre de Lie, $x\longrightarrow x_M$ une représentation de $g$ de dimension finie, $(x_1,\ldots,x_n)$ une base de $M$, et $(x_1^*,\ldots,x_n^*)$ la base duale.

a) Si $f$ est une forme linéaire sur $\overset{r}{\bigotimes}M$, on a :
\[
f=\sum_{1\leq i_1,\ldots,i_r\leq n}f(x_{i_1}\otimes\cdots\otimes x_{i_r})(x_{i_1}^*\otimes\cdots\otimes x_{i_r}^*).
\]

b) Soit $\beta$ une forme bilinéaire non dégénérée sur $M$, invariante pour $g$. Soit $(x'_1,\ldots,x'_n)$ la base de $M$ telle que $\beta(x'_i,x'_j)=\delta_{ij}$. Soit $f$ une forme linéaire invariante sur $\overset{r}{\bigotimes}M$. Déduire de a) que
\[
\sum_{1\leq i_1,\ldots,i_r\leq n}f(x_{i_1}\otimes\cdots\otimes x_{i_r})(x'_{i_1}\otimes\cdots\otimes x'_{i_r})
\]
est un élément invariant de $\overset{r}{\bigotimes}M$ indépendant du choix de la base $(x_i)$.

c) Soit $U$ l'algèbre enveloppante de $g$, $U^*$ l'espace dual de l'espace $U$. La représentation adjointe de $g$ se prolonge en une représentation $x\longrightarrow x_U$ de $g$ dans $U$ définie par $x_Uu=xu-ux$ pour tout $u\in U$. Donc $U^*$ est muni d'une structure de $g$-module. Pour qu'un élément $f$ de $U^*$ soit invariant, il faut et il suffit que $f(uv)=f(vu)$ quels que soient $u,v\in U$.
